% !TeX encoding = UTF-8
\documentclass[variant=guia,cover=institutional,mode=screen]{ua-engineering-v6-article}
% !TeX encoding = UTF-8
\UASetUniversity{Universidad Austral}
\UASetFaculty{Facultad de Ingeniería}
\UASetDepartment{Departamento de Computación}
\UASetCampus{Pilar}
\UASetCareer{Ingeniería Informática}
\UASetCommission{Comisión A}
\UASetProfessor{Equipo docente}
\UASetSemester{Segundo cuatrimestre 2026}
\UASetCourse{Sistemas Distribuidos}
\UASetDate{2026}


\UASetClassNumber{13}
\UASetClassTitle{Confiabilidad operativa}
\UASetDocKind{Guía resuelta}

\title{Clase 13 --- Guía resuelta}
\author{\UAProfessor}
\date{\UADate}

\begin{document}

\section{Criterio de lectura}
Las resoluciones son modelos de razonamiento. Deben verificarse contra los supuestos del caso y no memorizarse como recetas independientes del dominio.

\section{Ejercicios y resoluciones completas}
\begin{uaexercise}
Defina SLI, SLO y SLA, y explique la diferencia entre los tres conceptos.
\end{uaexercise}
\begin{uaanswer}
La respuesta debe situarse en SLI, SLO, SLA, error budgets, burn rates y operación guiada por objetivos. Se comienza declarando el supuesto decisivo y el comportamiento observable; luego se recorre una ejecución nominal y otra bajo SLI mal elegido, ventana incorrecta, budget consumido y alerta ruidosa. El mecanismo elegido debe vincularse con una garantía concreta, evitando afirmar propiedades fuera de su alcance.

La comparación debe usar al menos semántica, latencia, disponibilidad, complejidad operativa y recuperación. La alternativa preferible cambia cuando cambia el costo del error en el dominio; no existe ganador universal.

La evidencia apropiada sería SLI calculado, SLO, budget, burn rate y política de decisión. También debe indicarse una condición que refute la elección o haga preferible una solución más simple.
\end{uaanswer}

\begin{uaexercise}
Explique por qué un SLI debe elegirse desde la perspectiva del usuario y no solo desde la infraestructura.
\end{uaexercise}
\begin{uaanswer}
La respuesta debe situarse en SLI, SLO, SLA, error budgets, burn rates y operación guiada por objetivos. Se comienza declarando el supuesto decisivo y el comportamiento observable; luego se recorre una ejecución nominal y otra bajo SLI mal elegido, ventana incorrecta, budget consumido y alerta ruidosa. El mecanismo elegido debe vincularse con una garantía concreta, evitando afirmar propiedades fuera de su alcance.

El argumento debe identificar la hipótesis, construir la cadena lógica o el contraejemplo y concluir exactamente qué propiedad queda demostrada. Una intuición sin secuencia verificable no es suficiente.

La evidencia apropiada sería SLI calculado, SLO, budget, burn rate y política de decisión. También debe indicarse una condición que refute la elección o haga preferible una solución más simple.
\end{uaanswer}

\begin{uaexercise}
Dé dos ejemplos de buenos SLI y dos ejemplos de malos SLI.
\end{uaexercise}
\begin{uaanswer}
La respuesta debe situarse en SLI, SLO, SLA, error budgets, burn rates y operación guiada por objetivos. Se comienza declarando el supuesto decisivo y el comportamiento observable; luego se recorre una ejecución nominal y otra bajo SLI mal elegido, ventana incorrecta, budget consumido y alerta ruidosa. El mecanismo elegido debe vincularse con una garantía concreta, evitando afirmar propiedades fuera de su alcance.

La evidencia apropiada sería SLI calculado, SLO, budget, burn rate y política de decisión. También debe indicarse una condición que refute la elección o haga preferible una solución más simple.
\end{uaanswer}

\begin{uaexercise}
Explique por qué “CPU media del cluster” no siempre es un buen SLI.
\end{uaexercise}
\begin{uaanswer}
La idempotencia separa intentos físicos de una operación lógica. Se usa una clave estable, se persiste el resultado junto con el efecto y los reintentos devuelven ese mismo resultado. El costo es estado adicional, política de expiración y coordinación con el almacenamiento.

El argumento debe identificar la hipótesis, construir la cadena lógica o el contraejemplo y concluir exactamente qué propiedad queda demostrada. Una intuición sin secuencia verificable no es suficiente.

La evidencia apropiada sería SLI calculado, SLO, budget, burn rate y política de decisión. También debe indicarse una condición que refute la elección o haga preferible una solución más simple.
\end{uaanswer}

\begin{uaexercise}
Diseñe un SLO razonable para una API pública de lectura.
\end{uaexercise}
\begin{uaanswer}
Linearizabilidad exige que cada operación parezca ocurrir en un punto entre invocación y respuesta y que el orden respete tiempo real. Para validar una historia se propone un orden secuencial compatible; si ninguna ubicación de las operaciones satisface ambos requisitos, la historia no es linearizable.

Una solución completa organiza la decisión con P-F-G-S-T: define el problema, enumera fallas, fija la garantía, presenta el protocolo y reconoce el trade-off. Debe incluir estados intermedios y qué ve el usuario si la operación no termina.

La evidencia apropiada sería SLI calculado, SLO, budget, burn rate y política de decisión. También debe indicarse una condición que refute la elección o haga preferible una solución más simple.
\end{uaanswer}

\begin{uaexercise}
Diseñe un SLO razonable para un sistema de checkout.
\end{uaexercise}
\begin{uaanswer}
La idempotencia separa intentos físicos de una operación lógica. Se usa una clave estable, se persiste el resultado junto con el efecto y los reintentos devuelven ese mismo resultado. El costo es estado adicional, política de expiración y coordinación con el almacenamiento.

Una solución completa organiza la decisión con P-F-G-S-T: define el problema, enumera fallas, fija la garantía, presenta el protocolo y reconoce el trade-off. Debe incluir estados intermedios y qué ve el usuario si la operación no termina.

La evidencia apropiada sería SLI calculado, SLO, budget, burn rate y política de decisión. También debe indicarse una condición que refute la elección o haga preferible una solución más simple.
\end{uaanswer}

\begin{uaexercise}
Explique por qué un SLO de 100\% suele ser una mala idea práctica.
\end{uaexercise}
\begin{uaanswer}
La idempotencia separa intentos físicos de una operación lógica. Se usa una clave estable, se persiste el resultado junto con el efecto y los reintentos devuelven ese mismo resultado. El costo es estado adicional, política de expiración y coordinación con el almacenamiento.

El argumento debe identificar la hipótesis, construir la cadena lógica o el contraejemplo y concluir exactamente qué propiedad queda demostrada. Una intuición sin secuencia verificable no es suficiente.

La evidencia apropiada sería SLI calculado, SLO, budget, burn rate y política de decisión. También debe indicarse una condición que refute la elección o haga preferible una solución más simple.
\end{uaanswer}

\begin{uaexercise}
Explique por qué elegir una ventana temporal adecuada es parte del diseño del SLO.
\end{uaexercise}
\begin{uaanswer}
La idempotencia separa intentos físicos de una operación lógica. Se usa una clave estable, se persiste el resultado junto con el efecto y los reintentos devuelven ese mismo resultado. El costo es estado adicional, política de expiración y coordinación con el almacenamiento.

Una solución completa organiza la decisión con P-F-G-S-T: define el problema, enumera fallas, fija la garantía, presenta el protocolo y reconoce el trade-off. Debe incluir estados intermedios y qué ve el usuario si la operación no termina.

El argumento debe identificar la hipótesis, construir la cadena lógica o el contraejemplo y concluir exactamente qué propiedad queda demostrada. Una intuición sin secuencia verificable no es suficiente.

La evidencia apropiada sería SLI calculado, SLO, budget, burn rate y política de decisión. También debe indicarse una condición que refute la elección o haga preferible una solución más simple.
\end{uaanswer}

\begin{uaexercise}
Defina error budget y explique su significado operativo.
\end{uaexercise}
\begin{uaanswer}
El error budget es la porción de eventos o tiempo que puede incumplir el SLO en la ventana. Su valor operativo aparece al decidir ritmo de cambio: un burn rate alto obliga a proteger confiabilidad; budget saludable permite asumir riesgo controlado.

La evidencia apropiada sería SLI calculado, SLO, budget, burn rate y política de decisión. También debe indicarse una condición que refute la elección o haga preferible una solución más simple.
\end{uaanswer}

\begin{uaexercise}
Calcule el error budget mensual para un servicio con SLO de 99.9\%.
\end{uaexercise}
\begin{uaanswer}
La idempotencia separa intentos físicos de una operación lógica. Se usa una clave estable, se persiste el resultado junto con el efecto y los reintentos devuelven ese mismo resultado. El costo es estado adicional, política de expiración y coordinación con el almacenamiento.

La evidencia apropiada sería SLI calculado, SLO, budget, burn rate y política de decisión. También debe indicarse una condición que refute la elección o haga preferible una solución más simple.
\end{uaanswer}

\begin{uaexercise}
Explique qué decisiones podría tomar un equipo si está consumiendo demasiado rápido su error budget.
\end{uaexercise}
\begin{uaanswer}
La idempotencia separa intentos físicos de una operación lógica. Se usa una clave estable, se persiste el resultado junto con el efecto y los reintentos devuelven ese mismo resultado. El costo es estado adicional, política de expiración y coordinación con el almacenamiento.

La evidencia apropiada sería SLI calculado, SLO, budget, burn rate y política de decisión. También debe indicarse una condición que refute la elección o haga preferible una solución más simple.
\end{uaanswer}

\begin{uaexercise}
Explique por qué el error budget conecta confiabilidad con velocidad de cambio.
\end{uaexercise}
\begin{uaanswer}
La idempotencia separa intentos físicos de una operación lógica. Se usa una clave estable, se persiste el resultado junto con el efecto y los reintentos devuelven ese mismo resultado. El costo es estado adicional, política de expiración y coordinación con el almacenamiento.

El argumento debe identificar la hipótesis, construir la cadena lógica o el contraejemplo y concluir exactamente qué propiedad queda demostrada. Una intuición sin secuencia verificable no es suficiente.

La evidencia apropiada sería SLI calculado, SLO, budget, burn rate y política de decisión. También debe indicarse una condición que refute la elección o haga preferible una solución más simple.
\end{uaanswer}

\begin{uaexercise}
Explique por qué un SLO basado en promedio de latencia puede ser engañoso.
\end{uaexercise}
\begin{uaanswer}
La idempotencia separa intentos físicos de una operación lógica. Se usa una clave estable, se persiste el resultado junto con el efecto y los reintentos devuelven ese mismo resultado. El costo es estado adicional, política de expiración y coordinación con el almacenamiento.

El argumento debe identificar la hipótesis, construir la cadena lógica o el contraejemplo y concluir exactamente qué propiedad queda demostrada. Una intuición sin secuencia verificable no es suficiente.

La evidencia apropiada sería SLI calculado, SLO, budget, burn rate y política de decisión. También debe indicarse una condición que refute la elección o haga preferible una solución más simple.
\end{uaanswer}

\begin{uaexercise}
Diseñe un SLI/SLO usando p95 o p99 para un servicio de pagos.
\end{uaexercise}
\begin{uaanswer}
La respuesta debe situarse en SLI, SLO, SLA, error budgets, burn rates y operación guiada por objetivos. Se comienza declarando el supuesto decisivo y el comportamiento observable; luego se recorre una ejecución nominal y otra bajo SLI mal elegido, ventana incorrecta, budget consumido y alerta ruidosa. El mecanismo elegido debe vincularse con una garantía concreta, evitando afirmar propiedades fuera de su alcance.

Una solución completa organiza la decisión con P-F-G-S-T: define el problema, enumera fallas, fija la garantía, presenta el protocolo y reconoce el trade-off. Debe incluir estados intermedios y qué ve el usuario si la operación no termina.

La evidencia apropiada sería SLI calculado, SLO, budget, burn rate y política de decisión. También debe indicarse una condición que refute la elección o haga preferible una solución más simple.
\end{uaanswer}

\begin{uaexercise}
Explique qué problema aparece si el equipo mide algo distinto de lo que el usuario experimenta.
\end{uaexercise}
\begin{uaanswer}
La idempotencia separa intentos físicos de una operación lógica. Se usa una clave estable, se persiste el resultado junto con el efecto y los reintentos devuelven ese mismo resultado. El costo es estado adicional, política de expiración y coordinación con el almacenamiento.

La evidencia apropiada sería SLI calculado, SLO, budget, burn rate y política de decisión. También debe indicarse una condición que refute la elección o haga preferible una solución más simple.
\end{uaanswer}

\begin{uaexercise}
Diseñe una estrategia mínima para medir correctamente un SLI de disponibilidad.
\end{uaexercise}
\begin{uaanswer}
La idempotencia separa intentos físicos de una operación lógica. Se usa una clave estable, se persiste el resultado junto con el efecto y los reintentos devuelven ese mismo resultado. El costo es estado adicional, política de expiración y coordinación con el almacenamiento.

Una solución completa organiza la decisión con P-F-G-S-T: define el problema, enumera fallas, fija la garantía, presenta el protocolo y reconoce el trade-off. Debe incluir estados intermedios y qué ve el usuario si la operación no termina.

La evidencia apropiada sería SLI calculado, SLO, budget, burn rate y política de decisión. También debe indicarse una condición que refute la elección o haga preferible una solución más simple.
\end{uaanswer}

\begin{uaexercise}
Explique qué significa operar un sistema guiado por SLOs.
\end{uaexercise}
\begin{uaanswer}
La idempotencia separa intentos físicos de una operación lógica. Se usa una clave estable, se persiste el resultado junto con el efecto y los reintentos devuelven ese mismo resultado. El costo es estado adicional, política de expiración y coordinación con el almacenamiento.

La evidencia apropiada sería SLI calculado, SLO, budget, burn rate y política de decisión. También debe indicarse una condición que refute la elección o haga preferible una solución más simple.
\end{uaanswer}

\begin{uaexercise}
Construya un ejemplo donde un incidente técnico parezca grave, pero no implique violación del SLO.
\end{uaexercise}
\begin{uaanswer}
La idempotencia separa intentos físicos de una operación lógica. Se usa una clave estable, se persiste el resultado junto con el efecto y los reintentos devuelven ese mismo resultado. El costo es estado adicional, política de expiración y coordinación con el almacenamiento.

Una solución completa organiza la decisión con P-F-G-S-T: define el problema, enumera fallas, fija la garantía, presenta el protocolo y reconoce el trade-off. Debe incluir estados intermedios y qué ve el usuario si la operación no termina.

La evidencia apropiada sería SLI calculado, SLO, budget, burn rate y política de decisión. También debe indicarse una condición que refute la elección o haga preferible una solución más simple.
\end{uaanswer}

\begin{uaexercise}
Construya un ejemplo donde un problema “pequeño” técnicamente sí implique un daño real al servicio medido por SLO.
\end{uaexercise}
\begin{uaanswer}
La idempotencia separa intentos físicos de una operación lógica. Se usa una clave estable, se persiste el resultado junto con el efecto y los reintentos devuelven ese mismo resultado. El costo es estado adicional, política de expiración y coordinación con el almacenamiento.

Una solución completa organiza la decisión con P-F-G-S-T: define el problema, enumera fallas, fija la garantía, presenta el protocolo y reconoce el trade-off. Debe incluir estados intermedios y qué ve el usuario si la operación no termina.

La evidencia apropiada sería SLI calculado, SLO, budget, burn rate y política de decisión. También debe indicarse una condición que refute la elección o haga preferible una solución más simple.
\end{uaanswer}

\begin{uaexercise}
Explique cómo SLOs ayudan a resolver discusiones entre producto e ingeniería.
\end{uaexercise}
\begin{uaanswer}
La idempotencia separa intentos físicos de una operación lógica. Se usa una clave estable, se persiste el resultado junto con el efecto y los reintentos devuelven ese mismo resultado. El costo es estado adicional, política de expiración y coordinación con el almacenamiento.

La evidencia apropiada sería SLI calculado, SLO, budget, burn rate y política de decisión. También debe indicarse una condición que refute la elección o haga preferible una solución más simple.
\end{uaanswer}

\end{document}
